\documentclass[12pt,a4paper]{article}
\usepackage{geometry}
\geometry{left=2cm,right=2cm,top=2cm,bottom=2cm}
\usepackage{amsmath,amssymb,amsfonts}
\usepackage{booktabs}
\usepackage{array}
\usepackage{hyperref}
\usepackage{url}
\usepackage{enumitem}
\usepackage{float}
\usepackage{fontspec}
\usepackage{tikz}
\usepackage{graphicx}
\usepackage{caption}
\usepackage{subcaption}

% ========= 日本語フォント (Windows標準: MS Gothic) =========
\usepackage{luatexja}
\usepackage{luatexja-fontspec}

\setmainfont{MS Gothic}[
Script=CJK,
Language=Japanese
]
\setsansfont{MS Gothic}[
Script=CJK,
Language=Japanese
]
\setmonofont{MS Gothic}[
Script=CJK,
Language=Japanese
]

\setmainjfont{MS Gothic}[
Script=CJK,
Language=Japanese
]
\setsansjfont{MS Gothic}[
Script=CJK,
Language=Japanese
]
\setmonojfont{MS Gothic}[
Script=CJK,
Language=Japanese
]

% ========= YUCT マクロ =========
\newcommand{\Ke}{K_{\mathrm{eff}}}
\newcommand{\betaf}{\beta}
\newcommand{\kappac}{\kappa_c}
\newcommand{\Sodd}{S_{\mathrm{odd}}}
\newcommand{\Seven}{S_{\mathrm{even}}}

\hypersetup{
	colorlinks=true,
	linkcolor=blue,
	citecolor=blue,
	urlcolor=blue,
}

\setlength{\parindent}{0pt}
\setlength{\parskip}{1ex}
\raggedbottom

\begin{document}
	
	\begin{center}
		{\Large \textbf{YUCT 普遍スケーリング法則の直接実験的検証}}
	\end{center}
	\begin{center}
		{\Large \textbf{\(\varepsilon = \kappa_c \alpha \Ke^{-\beta}\)（\(\beta = 0.67\)）：五つの独立した方法}}
	\end{center}
	
	\vspace{0.5cm}
	\noindent
	A. V. Yakushev\\
	Yakushev Research, \url{https://yuct.org/}\\
	\vspace{0.3cm}
	\textbf{DOI:} \url{https://doi.org/10.5281/zenodo.18444598}\\
	\noindent
	2026年6月（改訂版）
	
	\vspace{0.5cm}
	
	\begin{abstract}
		\noindent
		本ノートは、Yakushev 統一調整理論 (YUCT) の中心的予測である普遍的誤差スケーリング指数 \(\beta = 0.67\) を直接確認できる、五つの独立した数値検証済み実験手法を提示する。各手法は公開データのみを使用し、電卓のみを必要とし、学部学生が半日で再現可能である。検証は素粒子物理学、生物学、進化遺伝学、材料科学にわたる。
	\end{abstract}
	
	\vspace{0.5cm}
	
	\noindent
	\textbf{キーワード:} YUCT, 普遍スケーリング, \(\beta = 0.67\), 実験手法, 質量梯子, アロメトリー, 突然変異率, 陽子質量, 融点。
	
	\section{序論}
	
	Yakushev 統一調整理論 (YUCT) は普遍誤差法則
	\begin{equation}\label{eq:error-law}
		\varepsilon = \kappa_c \, \alpha \, \Ke^{-\beta},
	\end{equation}
	ここで \(\beta = 2/3 \approx 0.667\)、\(\kappa_c = 1/3\) に基づく。以下の五つの方法により、理論的導出に依存せず、あらゆる科学者がこの予測を直接検証できる。
	
	\section{方法1：フェルミオン質量梯子}
	\label{sec:massladder}
	
	\textbf{所要時間:} 30分\\
	\textbf{データソース:} Particle Data Group (PDG) 2024 \cite{PDG2024}\\
	\textbf{原理:} YUCT の代数的ループが基本スケーリング量子を定める
	\[
	q = \left( \frac{\Sodd}{\Seven} \right)^{1/3} = \left( \frac{3}{2} \right)^{1/3} \approx 1.1447.
	\]
	全ての荷電フェルミオンの質量は \(m_f = m_e \cdot q^{N_f}\) と量子化され、\(N_f \in \frac{1}{2}\mathbb{Z}\) である。
	
	\begin{table}[H]
		\centering
		\caption{フェルミオン質量と半整数指数}
		\label{tab:fermions}
		\begin{tabular}{l c c c c}
			\toprule
			フェルミオン & 質量 (GeV) & \(\ln(m_f/m_e)\) & \(N_f\) (計算値) & 最近の \(\frac{1}{2}\mathbb{Z}\) \\
			\midrule
			\(e\)   & \(5.11\times10^{-4}\) & 0.000  & 0.00  & 0 \\
			\(\mu\)  & 0.10566 & 5.329  & 39.43 & 39.5 \\
			\(\tau\) & 1.7768  & 8.153  & 60.31 & 60.0 \\
			\(u\)   & \(2.16\times10^{-3}\) & 1.439  & 10.65 & 10.5 \\
			\(d\)   & \(4.67\times10^{-3}\) & 2.209  & 16.34 & 16.5 \\
			\(s\)   & 0.0935  & 5.206  & 38.51 & 38.5 \\
			\(c\)   & 1.273   & 7.817  & 57.83 & 58.0 \\
			\(b\)   & 4.183   & 9.006  & 66.63 & 66.5 \\
			\(t\)   & 172.57  & 12.726 & 94.18 & 94.0 \\
			\bottomrule
		\end{tabular}
		\par\smallskip
		\footnotesize
		\(\ln q = \frac{1}{3}\ln(3/2) \approx 0.135155\)。\(N_f\) (計算値) は \(\ln(m_f/m_e)/\ln q\)。
	\end{table}
	
	\textbf{学生演習:}
	\begin{enumerate}[label=(\arabic*)]
		\item PDG 要約表で9つの荷電フェルミオン質量を調べる。
		\item 各粒子について \(N_f = \ln(m_f / m_e) / \ln q\) を計算する。
		\item 9つの値すべてが整数または半整数の \(\pm 0.5\) 以内にあることを確認する。
	\end{enumerate}
	
	\textbf{期待される結果:} 9つのフェルミオン全てが半整数梯子上に乗る。
	
	\section{方法2：アロメトリスケーリング（クライバーの法則）}
	\label{sec:kleiber}
	
	\textbf{所要時間:} 30分\\
	\textbf{データソース:} Kleiber (1947) \cite{Kleiber1947}\\
	\textbf{原理:} 代謝率 \(B\) は体重 \(M\) と \(B \propto M^{\beta d_f/2}\) の関係でスケールする。哺乳類では \(d_f \approx 2.2\) より \(b \approx 0.73\)、単細胞生物では \(d_f \approx 2\) より \(b \approx 0.67\) を得る。
	
	\begin{table}[H]
		\centering
		\caption{哺乳類に関するクライバーの原データ}
		\label{tab:kleiber}
		\begin{tabular}{l c c c}
			\toprule
			動物 & 体重 \(M\) (kg) & 代謝率 \(B\) (kcal/day) \\
			\midrule
			マウス       & 0.02   & 3.3 \\
			ラット         & 0.20   & 20 \\
			モルモット  & 0.65   & 48 \\
			ネコ         & 3.0    & 152 \\
			イヌ         & 15     & 500 \\
			ヒト       & 70     & 1700 \\
			ウマ       & 700    & 10000 \\
			\bottomrule
		\end{tabular}
	\end{table}
	
	\textbf{手順:}
	\begin{enumerate}[label=(\arabic*)]
		\item \(\ln B\) を \(\ln M\) に対してプロットする。傾きは \(0.74 \pm 0.02\)。
		\item \(\beta = 2b / d_f\) を抽出する。哺乳類の \(d_f = 2.2\) に対して、
		\(\beta = 2 \times 0.74 / 2.2 \approx 0.67\)。
	\end{enumerate}
	
	\section{方法3：脊椎動物の突然変異率}
	\label{sec:mutation}
	
	\textbf{所要時間:} 30分（データ解析のみ）\\
	\textbf{データソース:} YUCT 付録 T、\url{https://github.com/Alexey-Yakushev-YUCT/YPSDC} にて入手可能\\
	\textbf{原理:} 突然変異率 \(\mu\) は DNA 修復の調整効率に反比例する：\(\mu \propto K_{\text{repair}}^{-\beta}\)。
	
	\textbf{手順:}
	\begin{enumerate}[label=(\arabic*)]
		\item データセット \texttt{mutation\_rates.csv} をダウンロードする。
		\item \(\ln \mu\) を \(\ln K_{\text{repair}}\) に対してプロットする。
		\item 直線を当てはめる。
	\end{enumerate}
	
	\textbf{期待される結果:} 傾き \(= -0.67 \pm 0.05\)、\(R^2 = 0.89\)。
	
	\section{方法4：質量梯子からの陽子質量}
	\label{sec:proton}
	
	\textbf{所要時間:} 10分\\
	\textbf{データソース:} PDG 2024 \cite{PDG2024}\\
	\textbf{原理:} 質量梯子は複合粒子に拡張される。陽子質量 \(m_p\) は半整数指数 \(N_f = 55.5\) に対応する。
	
	\begin{table}[H]
		\centering
		\caption{陽子指数の計算}
		\label{tab:proton}
		\begin{tabular}{l c}
			\toprule
			量 & 値 \\
			\midrule
			陽子質量 \(m_p\) & 0.938272 GeV \\
			電子質量 \(m_e\) & \(5.11\times10^{-4}\) GeV \\
			\(\ln(m_p/m_e)\) & 7.515 \\
			\(\ln q = \frac{1}{3}\ln(3/2)\) & 0.135155 \\
			\(N_f = \ln(m_p/m_e)/\ln q\) & 55.61 \\
			最近の半整数 & 55.5 \\
			\bottomrule
		\end{tabular}
	\end{table}
	
	\textbf{手順:}
	\begin{enumerate}[label=(\arabic*)]
		\item 陽子質量と電子質量を調べる。
		\item \(N_f = \ln(m_p / m_e) / \ln q\) を計算する。
		\item 半整数 \(55.5\) からの差が \(0.11\) 以内であることを確認する。
	\end{enumerate}
	
	\section{方法5：単純金属の融点}
	\label{sec:melting}
	
	\textbf{所要時間:} 30分\\
	\textbf{データソース:} CRC 化学・物理学ハンドブック \cite{CRC}、YUCT 付録 C \cite{AppendixC}\\
	\textbf{原理:} YUCT では、相転移温度は調整効率と \(T_{\text{融解}} \propto \Ke^{\beta}\) の関係でスケールする。単純金属ではこのべき乗則がよく従う。
	
	\begin{table}[H]
		\centering
		\caption{10種類の単純金属の融解温度と調整効率。\(\Ke\) の値は付録 C より。}
		\label{tab:melting}
		\begin{tabular}{l c c c c}
			\toprule
			元素 & \(\Ke\) & \(T_{\text{融解}}\) (K) & \(\ln \Ke\) & \(\ln T_{\text{融解}}\) \\
			\midrule
			Li & 1.85 & 453.7 & 0.615 & 6.118 \\
			Na & 2.13 & 371.0 & 0.757 & 5.916 \\
			K  & 2.38 & 336.5 & 0.867 & 5.819 \\
			Mg & 2.57 & 923.0 & 0.943 & 6.827 \\
			Al & 3.01 & 933.5 & 1.102 & 6.839 \\
			Cu & 4.62 & 1358  & 1.530 & 7.214 \\
			Zn & 3.96 & 692.7 & 1.376 & 6.540 \\
			Fe & 4.26 & 1811  & 1.449 & 7.502 \\
			Ni & 4.50 & 1728  & 1.504 & 7.455 \\
			Au & 4.97 & 1337  & 1.604 & 7.199 \\
			\bottomrule
		\end{tabular}
	\end{table}
	
	\textbf{手順:}
	\begin{enumerate}[label=(\arabic*)]
		\item 表から \(\ln \Ke\) と \(\ln T_{\text{融解}}\) の値をコピーする。
		\item \(\ln T_{\text{融解}}\) を \(\ln \Ke\) に対してプロットする。
		\item 点を通る直線を当てはめる。
	\end{enumerate}
	
	\textbf{期待される結果:} 傾きは \(0.58 \pm 0.09\)（\(R^2 \approx 0.62\)）。
	理論値 \(\beta = 0.67\) は 95\% 信頼区間内に収まる。
	
	\section{結論}
	
	ここで提示した五つの方法は数値的に検証済みであり、公開データのみを必要とする。それらは同じ背後にある指数 \(\beta = 0.67\) が、素粒子の質量、哺乳類の代謝率、DNA修復の忠実度、陽子の質量、金属の融点を組織していることを示している。いかなる学生も半日でこれらの結果を再現できる。
	
	\section*{データの入手可能性}
	全てのデータセットとスクリプトは YPSDC リポジトリで入手可能\\
	\url{https://github.com/Alexey-Yakushev-YUCT/YPSDC}
	
	\begin{thebibliography}{9}
		\bibitem{PDG2024} Particle Data Group, \emph{Review of Particle Physics},
		Prog.\ Theor.\ Exp.\ Phys.\ \textbf{2024}, 083C01 (2024).
		\bibitem{Kleiber1947} M.~Kleiber, ``Body size and metabolic rate,''
		\emph{Physiol. Rev.} \textbf{27}, 511--541 (1947).
		\bibitem{AppendixT} A.\,V.\,Yakushev, ``Appendix T: The Fundamental Law
		of Evolution in YUCT,'' in \emph{YUCT}, Zenodo (2026).
		\bibitem{AppendixJ} A.\,V.\,Yakushev, ``Appendix J: Complete Derivation of
		Standard Model Flavor Parameters from YUCT Topological
		Offsets,'' in \emph{YUCT}, Zenodo (2026).
		\bibitem{CRC} CRC Handbook of Chemistry and Physics, 106th ed.,
		CRC Press (2025).
		\bibitem{AppendixC} A.\,V.\,Yakushev, ``Appendix C: The Coordination‑Based
		Periodic Table of Elements,'' in \emph{YUCT}, Zenodo (2026).
	\end{thebibliography}
	
\end{document}